\documentclass[10pt,twocolumn]{article}
\usepackage[utf8]{inputenc}
\usepackage{amsmath,amssymb,amsfonts}
% \usepackage{algorithmic}  % Optional: for algorithm pseudocode
\usepackage{graphicx}
\usepackage{textcomp}
\usepackage{xcolor}
\usepackage{geometry}
\usepackage{abstract}
\usepackage{url}
\usepackage{balance}

% IEEE 风格设置
\geometry{
    a4paper,
    margin=0.75in,
    top=0.75in,
    bottom=1in
}

% 标题格式
\makeatletter
\renewcommand{\@maketitle}{%
  \newpage
  \null
  \vskip 2em%
  \begin{center}%
  \let \footnote \thanks
    {\LARGE \@title \par}%
    \vskip 1.5em%
    {\large
      \lineskip .5em%
      \begin{tabular}[t]{c}%
        \@author
      \end{tabular}\par}%
    \vskip 1em%
  \end{center}%
  \par
  \vskip 1.5em}
\makeatother

% 摘要格式
\renewcommand{\abstractname}{Abstract}
\renewenvironment{abstract}{%
  \small
  \begin{center}%
    {\bfseries \abstractname\vspace{-.5em}\vspace{\z@}}%
  \end{center}%
  \quotation
}{%
  \endquotation
}

% 章节格式
\makeatletter
\renewcommand\section{\@startsection {section}{1}{\z@}%
                                   {-3.5ex \@plus -1ex \@minus -.2ex}%
                                   {2.3ex \@plus.2ex}%
                                   {\normalfont\large\bfseries}}
\renewcommand\subsection{\@startsection{subsection}{2}{\z@}%
                                     {-3.25ex\@plus -1ex \@minus -.2ex}%
                                     {1.5ex \@plus .2ex}%
                                     {\normalfont\normalsize\bfseries}}
\makeatother

% 参考文献格式
\makeatletter
\renewcommand\@biblabel[1]{[#1]}
\makeatother

\begin{document}

\title{Your Research Paper Title}

\author{
\IEEEauthorblockN{1\textsuperscript{st} Your Name}
\IEEEauthorblockA{\textit{Department} \\
\textit{University/Institution}\\
City, Country \\
email@example.com}
\and
\IEEEauthorblockN{2\textsuperscript{nd} Co-Author Name}
\IEEEauthorblockA{\textit{Department} \\
\textit{University/Institution}\\
City, Country \\
coauthor@example.com}
}

% 定义 IEEEauthorblockN 和 IEEEauthorblockA（如果不存在）
\providecommand{\IEEEauthorblockN}[1]{#1}
\providecommand{\IEEEauthorblockA}[1]{#1}

\maketitle

\begin{abstract}
% Write your abstract here. The abstract should provide a concise summary of your research, including the problem, methodology, key findings, and conclusions. Typically 150-250 words.
% IEEE typically requires abstracts to be between 150-250 words.
\end{abstract}

\textbf{\textit{Keywords---}}keyword1, keyword2, keyword3, keyword4, keyword5

\section{Introduction}
\IEEEPARstart{T}{his} section introduces the background, motivation, and contributions of your research.

\subsection{Background}
% Background information about the research area

\subsection{Motivation}
% Why is this research important? What problem does it address?

\subsection{Contributions}
% What are the main contributions of this work?

\subsection{Paper Organization}
% Overview of how the paper is organized

\section{Problem Formulation}
This section provides a detailed formulation of the problem being addressed.

\subsection{Problem Statement}
% Clear and precise statement of the problem

\subsection{Problem Constraints}
% Constraints and assumptions

\subsection{Mathematical Formulation}
% Mathematical models and formulations
% Example:
% \begin{equation}
% \label{eq:example}
% f(x) = \sum_{i=1}^{n} w_i x_i
% \end{equation}

\section{Proposed Method}
This section describes the proposed methodology in detail.

\subsection{Method Overview}
% High-level overview of your proposed method

\subsection{Detailed Methodology}
% Detailed description of each component

\subsection{Algorithm}
% Algorithm pseudocode or detailed steps
% Example algorithm:
% \begin{algorithm}
% \caption{Your Algorithm Name}
% \begin{algorithmic}[1]
% \STATE Initialize parameters
% \STATE \textbf{for} each iteration \textbf{do}
% \STATE \quad Perform computation
% \STATE \textbf{end for}
% \end{algorithmic}
% \end{algorithm}

\subsection{Key Innovations}
% What makes your method novel or better?

\section{Experimental Setup}
This section describes the experimental setup, datasets, and evaluation metrics.

\subsection{Dataset}
% Description of datasets used

\subsubsection{Dataset Description}
% Details about the dataset

\subsubsection{Data Preprocessing}
% Data preprocessing steps

\subsection{Implementation Details}
% Implementation environment, tools, libraries
% Example:
% \begin{itemize}
% \item Python 3.9
% \item PyTorch 1.12
% \item CUDA 11.6
% \end{itemize}

\subsection{Evaluation Metrics}
% Metrics used for evaluation (MAE, MAPE, RMSE, R², RRSE, etc.)
% Example:
% \begin{itemize}
% \item Mean Absolute Error (MAE)
% \item Mean Absolute Percentage Error (MAPE)
% \item Root Mean Squared Error (RMSE)
% \item Coefficient of Determination (R²)
% \item Root Relative Squared Error (RRSE)
% \end{itemize}

\subsection{Baseline Methods}
% Methods used for comparison
% Example:
% \begin{itemize}
% \item XGBoost
% \item Random Forest
% \item LightGBM
% \item Multi-Layer Perceptron (MLP)
% \item CatBoost
% \end{itemize}

\subsection{Experimental Configuration}
% Hyperparameters, training settings, etc.

\section{Experimental Results}
This section presents the experimental results and analysis.

\subsection{Overall Performance}
% Overall performance comparison
% Example table:
% \begin{table}[h]
% \centering
% \caption{Overall Performance Comparison}
% \label{tab:overall}
% \begin{tabular}{|c|c|c|c|c|}
% \hline
% Method & MAE & MAPE & RMSE & R² \\
% \hline
% Method 1 & & & & \\
% \hline
% Method 2 & & & & \\
% \hline
% \end{tabular}
% \end{table}

\subsection{Comparison with Baseline Methods}
% Detailed comparison with other methods

\subsection{Ablation Studies}
% Analysis of different components

\subsection{Case Studies}
% Specific examples or case studies

\subsection{Visualization}
% Figures and tables showing results
% Example figure:
% \begin{figure}[h]
% \centering
% \includegraphics[width=\columnwidth]{figure.png}
% \caption{Your figure caption}
% \label{fig:example}
% \end{figure}

\section{Discussion}
This section discusses the results, insights, limitations, and future work.

\subsection{Result Analysis}
% Analysis and interpretation of the results

\subsection{Key Insights}
% Important insights from the experiments

\subsection{Limitations}
% Limitations of the proposed method

\subsection{Future Work}
% Potential future research directions

\section{Conclusion}
% Write your conclusion here. Summarize the main contributions, findings, and implications of your research.

\section*{Acknowledgment}
% Optional acknowledgment section

The authors would like to thank...

\begin{thebibliography}{00}
\bibitem{b1} Author, A., and Author, B., ``Title of paper,'' \textit{Journal/Conference Name}, vol. X, no. Y, pp. 1--10, Year.
\bibitem{b2} Author, C., \textit{Book Title}, Publisher, Year.
\bibitem{b3} Author, D., ``Title of webpage,'' Website Name. [Online]. Available: URL
\end{thebibliography}

% 定义 IEEEPARstart 命令（用于段落首字母大写）
\makeatletter
\newcommand{\IEEEPARstart}[2]{#1#2}
\makeatother

% Optional appendix section
% \appendices
% \section{Additional Experimental Results}
% % Additional results, figures, or tables

\balance
\end{document}
